\chapter{Modello cinematico del robot}

Come accennato nella \cref{sec:sistema_movimentazione}, il veicolo è stato progettato per essere in grado di eseguire sia manovre comuni come avanzamenti in linea retta e curve con orientamento tangenziale all'arco di curva, sia manovre più complesse come rotazioni sul posto e traslazioni laterali. Per raggiungere questo obiettivo, il sistema di movimentazione del robot è stato progettato con due ruote motrici indipendenti e due ruote castor libere di girare attorno al proprio perno verticale, disposte diagonalmente rispetto al centro geometrico del veicolo.

In \cref{fig:modello_cinematico} mostriamo una rappresentazione schematica del robot con le sole ruote attive, gli angoli di sterzo e le velocità lineari, dove il suo centro geometrico coincide con l'origine degli assi. In \cref{fig:modello_cinematico_2} riportiamo anche il caso generale in cui il veicolo è ad una certa distanza $x_c$ e $y_c$ dall'origine del sistema di riferimento globale.
\begin{figure}[htbp]
    \centering
    \includegraphics[width=\textwidth]{4/schema_modello.pdf}
    \caption{Rappresentazione schematica del modello cinematico del robot.}
    \label{fig:modello_cinematico}
\end{figure}
\begin{figure}[htbp]
    \centering
    \includegraphics[width=\textwidth]{4/schema_modello_2.pdf}
    \caption{Rappresentazione schematica del modello cinematico del robot in posizione generica.}
    \label{fig:modello_cinematico_2}
\end{figure}

Per derivare le equazioni cinematiche del veicolo, facciamo le seguenti ipotesi semplificative:
\begin{itemize}
    \item Le ruote motrici sono in grado di ruotare attorno al proprio asse verticale per variare l'angolo di sterzo $\delta_i,\, i\in{f,r}$.
    \item Le ruote motrici sono in grado di variare la propria velocità angolare $\omega_i,\, i\in{f,r}$ in modo indipendente.
    \item Le ruote castor sono libere di ruotare attorno al proprio perno verticale senza attrito.
    \item Le ruote sono perfettamente aderenti al terreno senza slittamenti laterali.
    \item Il veicolo si muove su un piano orizzontale senza pendenze.
    \item Le dimensioni del veicolo sono note e costanti pari ad $L$ e $W$, con le ruote motrici montate nei punti $ (L/2, W/2) $ e $ (-L/2, -W/2) $ rispetto al centro geometrico.
    \item Il centro di massa del veicolo coincide con il centro geometrico.
    \item Le forze esterne agenti sul veicolo (es. vento, urti, ecc.) sono trascurabili.
    \item Le accelerazioni del veicolo sono moderate, permettendo di trascurare gli effetti dinamici complessi come l'inerzia delle ruote e le forze di Coriolis.
    \item Il veicolo è simmetrico rispetto agli assi longitudinale e trasversale.
    \item Le ruote motrici sono identiche tra loro, con lo stesso raggio $r$ e caratteristiche di attrito.
    \item Il sistema di riferimento globale è inerziale e fisso, mentre il sistema di riferimento locale al robot si muove con esso.
\end{itemize}

Tali ipotesi sono comuni nella letteratura sui robot mobili e permettono di semplificare notevolmente l'analisi del modello cinematico. In particolare, per un robot dotato di quattro ruote indipendenti, tipicamente si considerano le ruote sullo stesso asse (quindi le due anteriori e le due posteriori) come un'unica ruota virtuale per asse, ognuna con un proprio angolo di sterzo e velocità~\cite{tourajizadeh2020fourwheel}. Questo approccio riduce il numero di variabili da considerare e semplifica le equazioni cinematiche del veicolo.

Tuttavia, quando questa semplificazione viene applicata, è importante tenere conto delle possibili differenze di comportamento tra le ruote sullo stesso asse, specialmente durante manove di traslazione laterale (anche dette \textit{crab steering}) o delle rotazioni sul posto. In questi casi, i modelli spesso diventano singolari e richiedono di essere adattati per gestire correttamente tali situazioni.

Per questo motivo, sebbene lo studio iniziale del problema sia stato condotto sul modello descritto in~\cite{tourajizadeh2020fourwheel}, si è deciso di verificare la validità di un modello più generale che invece è ottenuto tramite la trattazione classica delle equazioni di vincolo per i robot mobili~\cite{deluca1998feedback}.

Vediamo quindi il modello a biciclo descritto da ~\cite{tourajizadeh2020fourwheel}, dove le equazioni cinematiche del veicolo sono date da:

\begin{equation}
    \label{eq:betamodel}
    \begin{cases}
        \dot{x}      & = v \cos(\theta+\beta)                                                   \\
        \dot{y}      & = v \sin(\theta+\beta)                                                   \\
        \dot{\theta} & = \frac{v\cos(\beta)}{l_f+l_r}\left(\tan(\delta_f)-\tan(\delta_r)\right)
    \end{cases}
\end{equation}

dove:
\begin{eqnarray}
    v     & = & \frac{v_f \cos(\delta_f) + v_r \cos(\delta_r)}{2\cos(\beta)} \label{eq:v_betamodel}\\
    \beta & = & \tan^{-1}\left(\frac{l_f \tan(\delta_r) + l_r \tan(\delta_f)}{l_f + l_r}\right) \label{eq:beta_betamodel}
\end{eqnarray}

Rappresentano rispettivamente velocità del centro del robot e angolo di \textit{side-slip}, ovvero, l'angolo tra la direzione di avanzamento del veicolo descritta dal vettore velocità e l'asse longitudinale del robot.

In queste equazioni, $l_f$ e $l_r$ rappresentano rispettivamente la distanza tra il centro del veicolo e l'asse anteriore e posteriore delle ruote motrici, mentre $v_f$ e $v_r$ sono le velocità lineari delle ruote motrici anteriori e posteriori, per cui chiaramente valgono:
\begin{equation}\label{eq:tan_vel}
    v_f = r \omega_f \quad , \quad v_r = r \omega_r
\end{equation}
dove $r$ è il raggio delle ruote motrici e $\omega_f$, $\omega_r$ sono le velocità angolari delle ruote motrici anteriori e posteriori.

Si può vedere poi che sostituendo le \cref{eq:v_betamodel,eq:beta_betamodel} nella \cref{eq:betamodel}, si può ottenere un sistema matriciale del tipo $\dot{q} = J(q,u)u$, per mettere in relazione le velocità lineare e angolare del robot ($\dot(q)$) con le velocità angolari delle ruote motrici e gli angoli di sterzo ($u$). In particolare, scrivendo $J(q,u)$ come:
\begin{equation*}
    J(q,u) = R(\theta)\hat{J}(u)
\end{equation*}
dove $R(\theta)$ è la matrice di rotazione tra il sistema di riferimento globale e quello locale al robot descritta da:
\begin{equation}\label{eq:Rtheta}
    R(\theta) = \begin{bmatrix}
        \cos(\theta) & -\sin(\theta) & 0 \\
        \sin(\theta) & \cos(\theta)  & 0 \\
        0            & 0             & 1
    \end{bmatrix}
\end{equation}
e $\hat{J}(u)$ è la matrice jacobiana del sistema nel sistema di riferimento locale al robot, data da:
\begin{equation}\label{eq:Jhat}
    \hat{J}(u) = \frac{r}{2}\begin{bmatrix}
        cos(\delta_f)                          & cos(\delta_r)                          \\
        \cos(\delta_f)\tan(\beta)              & \cos(\delta_r)\tan(\beta)              \\
        \frac{\cos(\delta_f)}{l}(\Delta_{tan}) & \frac{\cos(\delta_r)}{l}(\Delta_{tan})
    \end{bmatrix}
\end{equation}
dove abbiamo definito:
\begin{eqnarray*}
    l          & = & l_f + l_r \\
    \Delta_{tan} & = & \tan(\delta_f) - \tan(\delta_r)
\end{eqnarray*}
In questo modo, si ottiene una rappresentazione compatta del modello cinematico, che può essere utile per l'implementazione del controllo.
\begin{equation}
    \begin{bmatrix}
        \dot{x} \\
        \dot{y} \\
        \dot{\theta}
    \end{bmatrix} =
    \frac{r}{2}\begin{bmatrix}
        \left[\cos(\theta)-\sin(\theta)(\Sigma_{tan})\right]\cos(\delta_f) &
        \left[\cos(\theta)-\sin(\theta)(\Sigma_{tan})\right]\cos(\delta_r)                                        \\
        \left[\sin(\theta)+\cos(\theta)(\Sigma_{tan})\right]\cos(\delta_f) &
        \left[\sin(\theta)+\cos(\theta)(\Sigma_{tan})\right]\cos(\delta_r)                                        \\
        \frac{1}{l}\Delta_{tan}\cos(\delta_f)                              & \frac{1}{l}\Delta_{tan}\cos(\delta_r)
    \end{bmatrix}
    \begin{bmatrix}
        \omega_f \\
        \omega_r \\
    \end{bmatrix}
\end{equation}
con
\begin{eqnarray*}
    \Sigma_{tan} &= \frac{l_r \tan(\delta_f) + l_f \tan(\delta_r)}{l_f + l_r}
\end{eqnarray*}